\documentclass[11pt]{article}
\usepackage[T1]{fontenc}
\usepackage{lmodern,microtype}
\usepackage[letterpaper,margin=0.9in]{geometry}
\usepackage{amsmath,amssymb,amsthm,mathtools}
\usepackage{booktabs,longtable,array,enumitem}
\usepackage{xcolor,hyperref,fancyhdr}
\definecolor{navy}{HTML}{18324D}
\definecolor{linkblue}{HTML}{145A8D}
\hypersetup{colorlinks=true,linkcolor=linkblue,urlcolor=linkblue,citecolor=linkblue,
 pdftitle={Semantic Sign Recovery and Cofinal Bounds: Explicit Diophantine Forms of Con(ZFC)},
 pdfauthor={AI-assisted research prepared for Vladimir Reshetnikov}}
\pagestyle{fancy}\fancyhf{}
\fancyhead[L]{\small\textsc{Explicit Diophantine consistency equations}}
\fancyhead[R]{\small 24 October 2026}
\fancyfoot[C]{\thepage}
\setlength{\headheight}{14pt}
\setlength{\parskip}{4pt}
\setlength{\parindent}{0pt}
\setlist{itemsep=3pt,topsep=4pt}
\newtheorem{theorem}{Theorem}[section]
\newtheorem{lemma}[theorem]{Lemma}
\newtheorem{proposition}[theorem]{Proposition}
\newtheorem{corollary}[theorem]{Corollary}
\theoremstyle{remark}\newtheorem{remark}[theorem]{Remark}
\newcommand{\Z}{\mathbb Z}\newcommand{\N}{\mathbb N}
\newcommand{\Con}{\operatorname{Con}}
\newcommand{\ZFC}{\mathsf{ZFC}}
\newcommand{\ZF}{\mathsf{ZF}}
\newcommand{\bits}{\operatorname{bits}}
\newcommand{\vtt}[1]{\texttt{\detokenize{#1}}}
\newcommand{\filepath}[1]{\path{#1}}
\newcommand{\sqnorm}[2]{#1^{2}-(#2)}
\allowdisplaybreaks
\title{\color{navy}\bfseries Semantic Sign Recovery and Cofinal Bounds\\[5pt]
\large Explicit Diophantine forms of $\Con(\ZFC)$}
\author{Research prepared for Vladimir Reshetnikov and ProveIt\\
\small AI-assisted mathematical development, implementation, and independent checks}
\date{24 October 2026}

\begin{document}
\maketitle
\begin{abstract}
We give an explicit quartic polynomial in 286 unrestricted integer variables
whose absence of an integer zero is provably equivalent, in ZFC, to
$\Con(\ZFC)$. It is the sum of squares of 196 literal quadratic residuals;
the collected polynomial has 6,607 monomials and coefficient height 128,446.
Relative to the inspected 324-variable direct-calculus parent, it removes
38 variables and 12 residuals while preserving degree and coefficient height.
The essential change is to orient the Pell abscissas and then derive the
signs of encoded arguments from their existing integer power certificates.
An integral output at a base at least two rules out a negative exponent.
An improved cofinal bound makes a previously quadratic beta modulus linear.
We also supply a 235-natural-variable quartic, three ten-positive-variable
equations evaluated in 80, 81, or 83 binary arithmetic operations, and two
eleven-integer equations evaluated in 107 or 108 operations. These results
improve specified constructions, with the degree and support tradeoffs made
explicit. They do not establish a global simplicity optimum. The new
equivalence arguments are ordinary proofs formalizable in ZFC; this packet
does not contain a new end-to-end proof-assistant derivation.
\end{abstract}

\section{Results, domain, and the meaning of simplicity}

The principal deliverable is the fixed polynomial
\begin{equation}\label{eq:main}
 \boxed{\displaystyle P_{286}(z_0,\ldots,z_{285})
       =\sum_{j=0}^{195}r_j(z_0,\ldots,z_{285})^2.}
\end{equation}
Every $r_j$ is printed literally in Appendix~\ref{app:residuals}. Thus
\eqref{eq:main} is an ordinary polynomial with explicit integer coefficients,
not an exponential Diophantine notation or an unspecified compilation request.
Its only unknowns are the 286 displayed coordinates, all in $\Z$.

\begin{theorem}[Principal consistency equivalence]\label{thm:main}
With the syntactic finite-proof meaning of $\Con(\ZFC)$ used by the retained
ProveIt calibration,
\[
 \ZFC\vdash
 \left[(\forall \mathbf z\in\Z^{286})\;P_{286}(\mathbf z)\ne0\right]
 \longleftrightarrow\Con(\ZFC).
\]
The polynomial has exact total degree four, 6,607 collected nonzero monomials,
and maximum absolute coefficient 128,446.
\end{theorem}

The direct construction encodes the numerical proof calculus C222. It uses
neither a simulation of K215 nor Jones's universal Diophantine polynomial.
The compact companions use the repository's separate native calibration;
their much larger constants must be counted when assessing simplicity.

\begin{table}[ht]
\centering\small
\begin{tabular}{lrrrr}
\toprule
Direct equation and domain&Variables&Rows&Monomials&Height\\
\midrule
Inspected parent, $\Z$&324&208&5,808&128,446\\
Parity companion, $\Z$&319&207&5,779&128,446\\
\textbf{Principal, $\Z$}&\textbf{286}&\textbf{196}&\textbf{6,607}&\textbf{128,446}\\
Inspected minimum parent, $\N$&237&185&69,710&133,499,131\\
\textbf{Natural companion, $\N$}&\textbf{235}&\textbf{184}&\textbf{22,185}&\textbf{90,911,633}\\
\bottomrule
\end{tabular}
\caption{All rows are quadratic residuals and all final polynomials are
quartic. Here $\N=\{0,1,2,\ldots\}$. The integer parent is package 131;
the natural parent is package 133, at the pinned repository commit.}
\end{table}

The principal reduction is $38/324\approx11.73\%$ in unrestricted integer
variables. Its monomial count increases by 799. The 319-variable companion
instead improves variables, residuals, and support simultaneously, with
unchanged degree and height. The natural companion improves every varying
metric in its matched row. These comparisons concern the identified direct
calculus architecture: the repository already contains quartics with far
fewer variables and enormous coefficient encodings.

\begin{table}[ht]
\centering\small
\begin{tabular}{lrrr}
\toprule
Domain and construction&Variables&Exact degree&Binary operations\\
\midrule
Package 133 product-square parent&10 positive&436&84\\
Free auxiliary coordinate&10 positive&436&\textbf{83}\\
Normalized, product transport&10 positive&560&\textbf{81}\\
Normalized, square transport&10 positive&750&\textbf{80}\\
Package 132 native parent&11 integers&924&109\\
Parity plus main-root saturation&11 integers&924&\textbf{107}\\
Package 132 short-index parent&11 integers&840&110\\
Parity plus main-root saturation&11 integers&840&\textbf{108}\\
\bottomrule
\end{tabular}
\caption{Every live binary addition, subtraction, and multiplication counts,
including squares and products by fixed constants. Constant leaves are fixed
numerals, not unknowns. The 80- and 81-operation points trade degree for cost.}
\end{table}

There is no metric-independent meaning of ``simplest'' here. We therefore
report the domain, dimension, degree, coefficient height, support, and
evaluation cost separately. No argument below proves global minimality or
historical priority. The contribution is the new combination and its
complete witnesses, equations, and verification records.

\section{Logical calibration and exact polynomial representation}

The direct source is calibrated by the two existing, unconditional interfaces
\begin{center}\small
\vtt{KEndpoint.c222_iff_zf_refutation}\\
\vtt{ConstructibleChoice.zf_refutation_iff_zfc_refutation}.
\end{center}
The first identifies derivability of falsity in C222 with a finite ZF
refutation. The second identifies finite ZF and ZFC refutations, through
the constructible-universe argument. The corresponding source declarations
and the numerical calculus definitions are retained in the archive. We do
not assume consistency, reflection, or the truth of all ZFC theorems in
order to compose these syntactic equivalences.

The inherited arithmetic construction encodes a finite derivation by a
uniform factorial/CRT scheme, digit equations, and Pell certificates for
the finitely many required powers. Its local polynomial permits exactly
the numerical C222 rule instances. The proof permits a rule multiset
initially; the retained Horn-flow argument extracts a finite derivation.
That arithmetic-to-calculus theorem, rather than a halting assumption,
is the interface used here. The included source ledger distinguishes it
from the new lemmas below.

Our task is to prove both implications
\[
 P_{286}(\mathbf z)=0\text{ for some }\mathbf z\in\Z^{286}
 \quad\longleftrightarrow\quad
 \text{a finite C222 refutation exists}.
\]
Negation and the two interfaces give Theorem~\ref{thm:main}. All new number
theory is elementary arithmetic formalizable in ZFC. This statement concerns
the equivalence; it does not assert either side of the resulting consistency
sentence.

The compiler turns each retained inequality $G\ge0$ into
\[
 G=u^2+v^2+w(w+1),\qquad u,v,w\in\Z.
\]
This form represents every nonnegative integer. Indeed, the three-square
theorem applies to $4G+1$; one square is odd and the other two are even,
giving the displayed form. Conversely every summand on its right is
nonnegative for integer arguments. The compiler then factors identities
exactly and eliminates only monic equations $z-Q=0$ whose substitution
leaves all residuals quadratic. Every eliminated coordinate has an explicit
integer restoration. The final centering is an integral translation.
These operations preserve existence of integer zeros. A sum of integer
squares vanishes exactly when each residual vanishes.

\section{A cofinal bound that removes a multiplication}\label{sec:cofinal}

In the parent construction the horizon is restricted to
$t=16+h_0^2$. A finite C222 derivation can first be unfolded to a finite
tree; $h_0$ is then chosen so that $t$ exceeds its step count, every local
schema parameter, and every full constructor label. Truth steps pad the
tree to exactly $t$ rows. This is an existence-preserving cofinal
restriction, with the witnesses chosen in the stated order.

Set
\[
 F=t!,\qquad V=(F^2+1)^t.
\]
The parent used the beta modulus $y=V(V-1)$. The following stronger estimate
permits the linear expression $y=V-1$.

\begin{lemma}[Linear cofinal modulus]\label{lem:linear}
Every pending multiset in the padded construction has an accumulator
strictly less than $V-1$. Also $F\mid V-1$ and $V-1\ge F+1$.
\end{lemma}
\begin{proof}
There are at most $t+1$ pending members, with multiplicity. Each member
has label at most $t$, so its weight is at most $1+tF\le(t+1)F$.
For $t\ge16$, $F>(t+1)^2$. This follows at $t=16$ directly and propagates
by multiplication by $t+1$. Hence
\[
 ((t+1)F)^{t+1}<F^{2t}\le(F^2+1)^t-1=V-1.
\]
For the strict inequality, after dividing by $F^{t+1}$ it suffices that
$F^{t-1}>(t+1)^{t+1}$; the stronger bound
$(t+1)^{2t-2}\ge(t+1)^{t+1}$ holds for $t\ge3$.
The last weak inequality follows by expansion. Finally $V\equiv1\pmod F$,
and the stated lower bound is immediate.
\end{proof}

The beta moduli $1+(i+1)y$, $0\le i\le t$, remain pairwise coprime:
a common prime divides a difference of indices, hence is at most $t$,
and divides both $F$ and $y$, contradicting divisibility of a modulus.
Thus exactly the original beta CRT construction works with $y=V-1$.
The resulting affine relation allows one more coordinate to be eliminated.
This improvement applies to both new direct branches.

There is no circular use of $F=t!$ in the soundness direction. Before
factorial recovery, the raw power and ceiling rows already imply $F\ge1$,
and then $(F^2+1)^t-1\ge F+1$. The stronger factorial bound is needed only
when constructing witnesses from an actual refutation.

\section{An even frame and an odd anchor quotient}\label{sec:parity}

For an integer $b\ge2$ define the Pell sequences by
\[
 (b+\sqrt{b^2-1})^j=X_j(b)+Y_j(b)\sqrt{b^2-1},\qquad j\ge0.
\]
Their recurrences give
\[
 X_{2j}=2X_j^2-1,\qquad Y_{2j}=2X_jY_j.
\]
When $b$ is odd, every $X_j$ is odd and $Y_j$ is odd precisely when
$j$ is odd. Doubling therefore proves
\begin{equation}\label{eq:valuation}
 v_2(Y_j(b))=v_2(j),\qquad j\ge1,\quad b\text{ odd}.
\end{equation}
We also use the elementary divisibility
\begin{equation}\label{eq:divisibility}
 c=Y_i(b)\quad\Longrightarrow\quad c^2\mid Y_{ic}(b).
\end{equation}
Expand $(X_i+c\sqrt{b^2-1})^c$: the first odd binomial term has a factor
$c^2$, and every subsequent odd term has a factor $c^3$.

Let the atoms be the parent's distinct actual power bases, exponents, and
outputs, including their affine restorations. We replace its frame by
\begin{align}
 W&=2+2\sum\text{atom}^2,\label{eq:frameW}\\
 A^2-W^3(W+2)(2z+1)^2&=1,\qquad A=2a^2-1.\label{eq:frameA}
\end{align}
The second equation replaces the old ternary guard on $A$ with one integer
$a$. Although $A=2a^2-1$ allows $A=-1$ or $1$ algebraically, the frame
norm excludes both; thus $A\ge2$ and is odd before any Pell classification.
The retained frame argument gives
\[
 W\ge514,\qquad A>W^W,\qquad A-1>4W^3.
\]

The shared anchor still has ordinate $C>W$ and cap
\begin{equation}\label{eq:cap}
 C=W+1+\sum_i y_i^2+u^2+v^2+w(w+1).
\end{equation}
Replace its guarded common multiple by
\begin{equation}\label{eq:oddJ}
 J=2C^2(2e+1),\qquad e\in\Z.
\end{equation}
It is nonzero and divisible by $C^2$, irrespective of its sign. The common
Pell norm and the old comparison norm use absolute indices and independent
coordinate signs. Their index comparison therefore still proves
$C=Y_W(A)$; it does not require $J>0$.

\begin{lemma}[Completeness of the parity restrictions]\label{lem:parity}
Every genuine front tuple has private witnesses satisfying
\eqref{eq:frameW}--\eqref{eq:oddJ} and the retained anchor comparisons.
\end{lemma}
\begin{proof}
Choose $A=X_W(W+1)$ and $a=X_{W/2}(W+1)$; $W$ is even, so Pell doubling
gives $A=2a^2-1$. The recurrence at $W+1$ gives
$Y_j(W+1)\equiv j\pmod{2W}$; hence $Y_W(W+1)/W$ is odd and supplies
the integer $z$ in \eqref{eq:frameA}.
Now take $C=Y_W(A)$, $\ell=2WC$, $J=Y_\ell(A)$, and $H=X_\ell(A)$.
Equation~\eqref{eq:divisibility} and doubling give $2C^2\mid J$.
Equation~\eqref{eq:valuation} gives
\[
 v_2(J)=1+v_2(W)+v_2(C)=1+2v_2(C).
\]
Thus $J/(2C^2)$ is positive odd and defines $e$. The abscissa $H$ is odd
and coprime to $J$. CRT supplies the old comparison center $D$ with
$D\equiv A\pmod H$ and $D\equiv1+4J\pmod{8J}$.
The parent's growth bound makes the cap slack nonnegative, and its
comparison pair at index $W$ supplies the remaining witnesses.
\end{proof}

These two compatible choices save five unrestricted integer coordinates
and one residual. By themselves they give the separate 319-variable
companion. Their role in the principal equation is to support the larger
semantic sign reduction next.

\section{Recovering signs from existing power equations}\label{sec:signs}

There are fifteen principal pairs $(x_i,y_i)$ with proposed signed indices
$k_i$. We retain every principal norm and power congruence, and strengthen
one already present pooled residue to
\begin{equation}\label{eq:pooled}
 \sum_i\bigl((y_i-k_i)^2+(x_i-1)^2\bigr)=(A-1)q_0,
 \qquad q_0\in\Z.
\end{equation}
Adding the abscissa squares introduces no coordinate. It does increase
the collected support, which is reported rather than suppressed.

\begin{lemma}[Orientation and exact signed indices]\label{lem:orientation}
At a zero of the new source,
$x_i=X_{|k_i|}(A)$ and $y_i=Y_{k_i}(A)$, where
$Y_{-j}=-Y_j$.
\end{lemma}
\begin{proof}
Pell classification gives
$x_i=\sigma_iX_{j_i}(A)$ and $y_i=\epsilon_iY_{j_i}(A)$,
with $j_i\ge0$ and signs $\sigma_i,\epsilon_i\in\{\pm1\}$.
The cap bounds $j_i<W$, while the frame gives $|k_i|<W$.
Modulo $A-1$, the Pell recurrences give $X_j\equiv1$ and $Y_j\equiv j$.
Thus $A-1$ divides
\[
 E_0=\sum_i\bigl((\epsilon_i j_i-k_i)^2+(\sigma_i-1)^2\bigr).
\]
Each summand is less than $4W^2$. Hence
$0\le E_0<60W^2<4W^3<A-1$, so $E_0=0$.
This forces every $\sigma_i=1$ and every signed index as claimed.
The huge actual abscissas need not be bounded: the argument bounds their
small residues. The zero index presents no exception.
\end{proof}

\begin{lemma}[Integer outputs exclude negative exponents]\label{lem:negative}
Let a retained power request have integer base $n\ge2$, signed index $k$,
and integer output $m$. Under the frame and
Lemma~\ref{lem:orientation}, it forces $k\ge0$ and $m=n^k$.
\end{lemma}
\begin{proof}
The frame gives $n,|k|,|m|<W$. Put $T_n=2An-n^2-1>W^W+W$.
For $k\ge0$, the Pell recurrence gives $m\equiv n^k\pmod{T_n}$;
the absolute difference is smaller than $T_n$, so equality follows.
For $k=-h<0$, use
\[
 \bigl(X_h-(A-n)Y_h\bigr)\bigl(X_h+(A-n)Y_h\bigr)
 =1+T_nY_h^2.
\]
The two factors reduce to $n^h$ and $m$, respectively. Therefore
$mn^h\equiv1\pmod{T_n}$. But
\[
 |mn^h-1|\le(W-1)^W+1<W^W<T_n.
\]
Consequently $mn^h=1$ as integers, impossible for $n\ge2$ and $h\ge1$.
For an already nonnegative index the first argument also works at base one.
\end{proof}

This lemma is the main compression mechanism. A separate ternary witness
is unnecessary when an existing integer power equation already proves the
corresponding exponent nonnegative.

\subsection{The acyclic bootstrap, including zero cases}

Write the raw front as
\[
 t=16+h_0^2,\quad B=r+s+1,\quad FB=V+s,\quad r,s\ge0,
\]
and retain the endpoint
\begin{equation}\label{eq:endpoint}
 x-1=(1+(t+1)y)u_t.
\end{equation}
The nine actual CRT representatives are $a_1,\ldots,a_9$; their stored
coordinates are $a_i-1$. Set
\begin{equation}\label{eq:radix}
 H_0=3+2\sum_{i=1}^9a_i,\quad R_0=2+t+x+y+F,\quad K=2R_0,
 \quad q=(BH_0)^K,\quad q=(2v+1)^2.
\end{equation}
The power in \eqref{eq:radix} names a retained graph, not an operation of
the final polynomial. The odd-square equation uses one new integer $v$
and replaces a ternary guard on $q-B-1$.

The following order avoids every dependency cycle.
\begin{enumerate}
\item The square gives $q\ge1$, while $t\ge16$ independently. The oriented
graph for $(F^2+1)^t$ gives $V\ge1$, even before $F$ has a sign. The ceiling
rows imply $F\ge1$ and $B>F\ge1$. With $y=V-1$ we obtain $y\ge F+1$.
\item Every actual representative $a_i$ is the exponent of a retained
request at base $q+1\ge2$. Lemma~\ref{lem:negative} gives $a_i\ge0$.
The exact source call indices are $5,7,8,9,10,12,13,15,16$; the verifier
checks their bases and affine exponent restorations against the literal
power ledger. Now $H_0$ is odd and at least three.
\item The graph for $B^K$, with $B\ge2$, gives $K\ge0$. If $K=0$, then
$x=-2-t-y-F$. Equation~\eqref{eq:endpoint} would make
$L=1+(t+1)y$ divide $x-1$, although
\[
 0<|x-1|=3+t+y+F<L,
\]
because $L-(3+t+y+F)=t(y-1)-F-2\ge(t-1)F-2>0$.
Hence $K\ge2$, since $K$ is even. Its second graph gives
$q=(BH_0)^K\ge36$.
\item Power call 14 has base $q$ and exponent $x+1$.
Lemma~\ref{lem:negative} gives $x\ge-1$. If $x=-1$, the endpoint would
make $L>2$ divide $-2$. Thus $x\ge0$. The removed output slack is also
nonnegative, since $q>B$.
\item The remaining constructor and elementary indices are now
nonnegative, and all eighteen power requests acquire their intended
meaning. Each paired remainder row has
$r_i+s_i=L_i-1\ge0$ with $L_i$ a positive power of $q$; its retained
guard $r_is_i\ge0$ forces both remainders nonnegative. Nonnegative-numerator division
rows restore the remaining quotient signs.
\end{enumerate}

Only fourteen semantic guards remain: the two initial ceiling guards,
eleven remainder products, and the shared cap. The two initial guards are
not removed by an argument that depends on their consequences.

\subsection{Why allowing a zero representative is sound}

The argument proves $a_i\ge0$, whereas the parent required $a_i\ge1$.
This is a deliberate extension, not a disguised identification of the
stored coordinate $-1$ with a natural number. The CRT theorem uses the
\emph{actual} argument $a_i$ in binomial coefficients and local constructor
polynomials. Its range lemma includes $a_i=0$; coefficients of positive
order then vanish. The local rule parameters in C222 range over natural
numbers including zero.

The choice $H_0=3+2\sum a_i$ provides the needed margin. For ordinary
range coefficients, $0\le a_i<H_0$ suffices. For the derived label, put
$n=\max(1,a_a,a_b,a_c,a_d)$. Then $n+1\le H_0$, including the case in
which all four arguments are zero. The parent's fixed constructor bound
$J_{\rm label}\le(n+1)^{4194304}$ and its factorial estimate
$4194304(t+1)<2F<K$ remain valid. Thus every required digit coefficient
fits below $H_0^K\le q$. The unchanged prime-power CRT and Horn-flow
arguments yield a C222 refutation.

The exact binomial and factorial recovery precedes that final decoding.
The ceiling gives $V\le BF<B^2$, while $q\ge B^2>V$. The uniform
coefficient bound therefore makes the retained exact digit graph give
$B=\binom{F^2+t}{t}$. The parent's padded factorial fixed-point theorem
then gives $F=t!$. This step uses neither positive representatives nor
their old guards.

\subsection{Completeness and the final count}

Start with a finite C222 refutation. Lemma~\ref{lem:linear} constructs the
new beta modulus and CRT code. Choose the independent representatives
strictly positive as in the parent; permitting zero in the reverse
direction imposes no new requirement here. Fix them before choosing
$H_0$ and $q$. The factorial product
\[
 B=\prod_{i=1}^t(1+F^2/i)
\]
is odd: $i\mid F$ and $F$ is even. Hence $BH_0$ is odd and
\[
 v=\frac{(BH_0)^{R_0}-1}{2}\in\Z
\]
supplies the odd-square radix. Lemma~\ref{lem:parity} supplies the private
frame and anchor. Canonical Pell pairs satisfy both congruences in
\eqref{eq:pooled}; the quotient and all remaining guards therefore exist.
This constructs a zero from every refutation and completes the new
semantic part of Theorem~\ref{thm:main}.

\begin{center}\small
\begin{tabular}{lrr}
\toprule
Change from the 324-variable parent&Variables saved&Rows saved\\
\midrule
Even frame and odd anchor quotient&5&1\\
Nine representative signs recovered from powers&27&9\\
The sign of $x$ recovered from its power and endpoint&3&1\\
Odd-square radix in place of a ternary guard&2&0\\
Linear cofinal modulus&1&1\\
\midrule
Total&38&12\\
\bottomrule
\end{tabular}
\end{center}

\section{The natural-variable companion}\label{sec:natural}
The natural branch uses the separate raw source of package 133, with its
dominant common digit modulus $M=Nq^{x+1}$. Its six exact powers and eleven modular numerator requests share fifteen
principal indices: four power indices, one exact-binomial upper index,
nine CRT representatives, and one constructor label. All raw coordinates in this
branch are nonnegative. The new linear beta modulus of
Section~\ref{sec:cofinal} is combined with a single bound-and-congruence
row; the one-row mechanism is credited to package 128.

Let $\mathcal S$ be the twelve distinct actual bases, exponents, and outputs
of the six power requests. Put
\[
 B_*=2+\sum\mathcal S+\sum_{i=1}^9a_i+J_{\rm label},
 \qquad W=2B_*.
\]
Every requested index $e_i$ is positive and less than $B_*$ directly from
the raw definitions. This assertion uses no power or factorial semantics.
The retained lower frame has
\[
 L^2-\bigl((W+1)^2-1\bigr)(WZ)^2=1,\qquad Z=z_0+1,
 \qquad A=L+\sigma+1,\quad \sigma\ge0.
\]
The shift by one is an existence-preserving restriction: the later common
root congruence admits cofinally many compatible $A$. The lower frame
gives $L\ge X_W(W+1)\ge(W+1)^W$ and $A-1\ge L>W^3$.

Retain the complete indexed anchor at $W$, with ordinate $C$, and replace
the old separate cap and the two staged residues by the single quadratic
row
\begin{equation}\label{eq:naturalcap}
 \boxed{C=W+\sum_{i=1}^{15}(Y_i-e_i)^2+(A-1)Q,\qquad Q\ge0.}
\end{equation}
The nonnegative domain is essential here. It first supplies $C\ge W$,
so the anchor theorem yields $C=Y_W(A)$ before any seed-index bound is
invoked.

\begin{lemma}[One-row identification of all indices]
Under these raw natural conditions, \eqref{eq:naturalcap} forces
$Y_i=Y_{e_i}(A)$ for every principal pair. Conversely, canonical pairs
admit a positive integer $Q$ satisfying that row.
\end{lemma}
\begin{proof}
Write $B=B_*$ temporarily. If $1\le e<B$, put $z=Y_B(A)$ and
$x=X_B(A)$. The norm gives $(A-1)z<x\le Az$. Since $z\ge B$,
\[
 Y_{B+1}(A)-e>(2A-2)z+1,\qquad Y_{2B}(A)=2xz\le2Az^2.
\]
For $A\ge2$, $(2A-2)^2\ge2A$, so
$(Y_{B+1}-e)^2>Y_{2B}=C$. But \eqref{eq:naturalcap} makes each squared
error smaller than $C$. Thus the actual index $j_i$ is at most $B$.
Reducing modulo $A-1$ now gives
\[
 A-1\mid\sum_i(j_i-e_i)^2,
 \qquad 0\le\sum_i(j_i-e_i)^2\le15(B-1)^2<W^3<A-1.
\]
Every error is zero.

For the converse choose all pairs canonically and set $C=Y_{2B}(A)$.
With $z=Y_{B-1}(A)$, the sum of the squared errors is at most $15z^2$,
whereas $C>2(A-1)A^2z^2$. Since $A>W+1$, their difference exceeds $W$.
Both $C-W$ and each $Y_{e_i}-e_i$ are divisible by $A-1$, so the
quotient in \eqref{eq:naturalcap} is a positive integer.
\end{proof}

All six powers are now exact by the retained small-residue identities;
the dominant common modulus, factorial fixed point, range digits, and
calculus decoding follow in the original order. In completeness, the
semantic data and all frame atoms are fixed first. The common-root
congruence then supplies a sufficiently large compatible $A$, followed
by the anchor and canonical seeds. There is no self-referential choice
of the frame.

For the linear beta restriction, use a new nonnegative coordinate
$v=V-1$ and the row $y=v$, replacing the old endpoint slack.
Soundness restores that slack as $y-F-1\ge0$ after the early ceiling
bootstrap; completeness is Lemma~\ref{lem:linear}. All occurrences of
$V$ and $A$ are shifted together in the literal graph.

Monic projection with explicit nonnegativity certificates gives the
235-variable equation with 184 residuals, 22,185 monomials, and height
90,911,633. A sparser 236-variable companion keeps two staged rows:
the first authenticates the five raw-bounded indices; only after exact
power and factorial recovery does the second authenticate the remaining
ten. It has 185 residuals, 15,645 monomials, and height 88,233,865.
Both have degree four and the same ZFC-provable consistency equivalence
over $\N$. Their complete literal residuals, expansions, inverse
projection certificates, and ordinary proofs are supplied in
\filepath{natural235/}.


\section{Three compact positive equations}\label{sec:compact}
This section is independent of the direct-calculus coefficients. It fixes
the entire indexed BFS943 tuple $(\beta,K_0,M_C,M_F)$ of package 133.
The finite recipe in Appendix~\ref{app:constants} and the supplied integer
payloads specify those four numbers completely. In particular
\[
 \beta\ge31\text{ is odd},\quad K_0>0,\quad
 0<M_C<\beta<M_F<2\beta.
\]
The largest numeral has 592,884 bits. The already shifted $M_F$ must not
be shifted a second time or mixed with another coefficient tuple.

\subsection{Complete formulas}

For the 83-operation equation the ten unknowns are
\[
 J,F,\alpha,i,T,w,Y,n,\eta,\zeta\in\Z_{>0}.
\]
For the normalized 80- and 81-operation equations replace the input $Y$
by $s\in\Z_{>0}$ and define $Y=q^3s$. The common polynomial definitions are
\begin{align*}
 q&=\beta J+1,& h&=q-1,& C&=q-F-\alpha,\\
 X&=wq,&\mu&=q^2+1,& k&=\eta+\zeta,\\
 c&=kY+\eta,& A&=\mu Y(X+1)+1,\\
 \Delta&=A^2-1,& H&=4A-5,\\
 v&=X\mu^2Y^2,&E&=1+v(v+1)k^2,\\
 R&=[h(q-F)+\alpha](q^2-1)+(M_C+qM_F)J,\\
 e&=k-R-1,&B&=(K_0+w)C-F,\\
 U&=ic^2,&a&=\Delta U^2E,&F_0&=1+a,\\
 G&=EF_0,\\
 F_*&=1+4aF_0,&Q&=\Delta(F_*-1),\\
 b&=cT-R,&V&=F_*b-c,&K_+&=X(3c-2X)+2.
\end{align*}
The abbreviation $F_*$ denotes a polynomial, not an additional unknown.
Take the transport weight $L=V^2$ for the 83- and 80-operation forms,
and $L=bV$ for the 81-operation form. For the unnormalized form set
\begin{align*}
 P&=q^3h,&M&=P\mu^2H,\\
 Z&=P\mu^2[H(n-LB)+K_+]
       +H[\mu^2(hY-q^3B)-Pe].
\end{align*}
For the normalized forms set
\[
 M=h\mu^2H,\qquad
 Z=h\mu^2[H(n-LB)+K_+]-H[\mu^2B+he].
\]
In each case the complete equation is
\begin{equation}\label{eq:compact}
 \boxed{Q\bigl(Z^2-M^2GV^2\bigr)+(M^2G-Z^2)=0.}
\end{equation}
It is identically $(Q-1)Z^2-M^2G(QV^2-1)=0$.
The strict positivity of the ten inputs is part of the domain.

\subsection{Root extraction and the free auxiliary coordinate}

The package-133 kernel uses Pell doubling:
\[
 F_*=(1+2a)^2,\qquad Q=4\Delta^2U^2G,\qquad\gcd(E,F_0)=1.
\]
Positive inputs give $q\ge32$, $c\ge3$, $a>c$, $Q>c^2$, and $M>0$.
If $b=0$, then $V=-c$; otherwise $|V|\ge F_*-c>1$. At a zero, put
$z=Z/M$. Since $G\mid Q$, both $Q-1$ and $QV^2-1$ are $-1$ modulo $G$.
Prime valuations in
\[
 (Q-1)z^2=G(QV^2-1)
\]
make $G=g^2$ for an integer $g>0$. Coprimality gives
$E=\tau^2$, $F_0=f_0^2$, and $g=\tau f_0$.

The positive rational $y=|z|/g$ satisfies
\begin{equation}\label{eq:adjacent}
 QV^2-(Q-1)y^2=1.
\end{equation}
It is integral. Indeed, if $y=r/m$ is reduced, then $m^2\mid Q-1$.
For $v_1=|V|>1$, the norm-preserving transformation
\[
 v_2=(2Q-1)v_1-2(Q-1)y,\quad y_2=(2Q-1)y-2Qv_1
\]
keeps $v_2$ integral, preserves the reduced denominator of $y_2$, and
has $-v_1<v_2<v_1$. The preserved norm excludes $v_2=0$. Descending in $|v_1|$ reaches one, where $y=1$.
Consequently $z=\epsilon gy$ for $\epsilon=\pm1$.

Define, in the proof only, $f=1+2a$ and $I=2Ug$. Then
\[
 f^2-\Delta I^2=1,\quad f^2=F_*,\quad Q=(\Delta I)^2,
 \quad c^2\mid I.
\]
The retained main-folding theorem permits either sign of $V$ and yields
$c=Y_p(A)$ for odd $p\ge1$. The first norm gives
$k=2Y_{n_0}(2v+1)$. The case $n_0=1$ would give $c=2Y+1$ strictly
between the first two Pell ordinates, so $n_0\ge2$. In particular
\[
 k\ge8v+4,\qquad c>k>8q^5,\qquad K_+>0.
\]

The new improvement is that $b=cT-R$ needs no preliminary shift
$T+\alpha$ when the transport argument is proved with absolute values.
At a zero, $b$ cannot vanish: otherwise
\[
 y^2=c^2+\frac{c^2-1}{Q-1}
\]
lies strictly between consecutive integers. For nonzero $b$, its sign
equals that of $V$ and $|b|<|V|$. Also
\[
 a^2-g^2=E[\Delta U^2E(\Delta U^2-1)-1]>0,
\]
so $g<a$ and $|V|\ge F_*-c>3g$. From \eqref{eq:adjacent},
$y<\tfrac32|V|$, hence $gy<V^2/2$.

Put $\kappa=K_+/H>0$, $t=B/h$, $\rho=e/\mu^2$, and
$W_*=1+hL$. The exact carrier identities are
\[
 \frac ZM=n+\kappa+Y/q^3-W_*t-\rho
 \quad\hbox{or}\quad
 \frac ZM=n+\kappa-W_*t-\rho.
\]
The one-sided estimate
\[
 b-e=cT-k+1>0,\qquad
 \rho<\frac b{\mu^2}\le\frac{|b|}{\mu^2}<\frac{|V|}{\mu^2}
\]
uses only $T\ge1$ and $c>k$. If $B\le-1$, then for $L=V^2$ either
carrier gives
\[
 n<\frac{V^2}{2}-V^2+\frac{|V|}{\mu^2}<0,
\]
a contradiction. Thus $B\ge0$, and the definition of $B$ forces $C>0$.
This proves the essential new soundness step for the 83- and 80-operation
equations. It never assumes the deleted shift has a positive inverse.

\subsection{Why the smaller weight works after normalization}

For $L=bV$ a stronger bound is needed. The literal identities $q\mid X,Y$
in the normalized chart exclude the exceptional main index $p=3$ before
transport. Suppose $p=3$, so $c=4A^2-1$. A first index $n_0\ge3$ would
give $kY>32v^2Y>4A^2>c$, so $n_0=2$ and $k=8v+4$.
Writing $D=(X+1)^2-2XY$ then gives the exact identity
\[
 \eta=4\mu^2Y^2D+8\mu Y(X+1)+3-4Y.
\]
The inequalities $0<\eta<k$ exclude $D\ge2X$ and $D\le0$.
For the latter, $D=1+X(X+2-2Y)\le-1$ implies $X+1<2Y$, making the
display negative. Thus $0<D<2X$, which forces
$2Y-X\in\{1,2\}$. This contradicts $q\mid X,Y$ and $q\ge32$.
Hence $p\ge5$ and $c>4\Delta+3$.

In \eqref{eq:adjacent}, put $S=\Delta I$ and $Q=S^2$. Pell classification
gives an odd auxiliary index; its first and third absolute $V$ values
are $1$ and $4Q-3$. The first is already excluded. The second is
impossible because $V=F_*b-c$ would give either
\[
 F_*(b-4\Delta)=c-4\Delta-3,
 \qquad F_*(b+4\Delta)=c+4\Delta+3,
\]
whose right side is strictly between zero and $F_*$. Therefore
$|V|\ge16Q^2-20Q+5>8Q^2$. As $c<F_*$,
\[
 |b|>8Q^2/F_*-1>4\Delta^3I^2-1>4g.
\]
Now $L=bV=|b||V|>0$, $gy<3L/8<L/2$, and
$\rho<|b|/\mu^2<L/2$. If $B\le-1$, the carrier again gives $n<0$.
This proves the missing sign step for the 81-operation equation.

\subsection{Later decoding, converse, and exact size}

After $C>0$, the retained mask bound gives $3q<R<q^4-2$. The signed-index
lemma uses the actual residues $V\equiv-c\pmod f$ and
$V\equiv-R\pmod c$, with $c^2\mid I$; it allows either sign of $V$ and
contains no condition $T>\alpha$. It gives $p=R$. The original parity,
denominator, first-index, and whole-product projector arguments then
restore the native source and its finite refutation. The normalized
carrier supplies a stronger denominator premise; no nonunit is cancelled
modulo a composite modulus.

Conversely the parent's canonical witnesses give a positive old
$T_{\rm old}$ and $q^3\mid Y$. Set $T=T_{\rm old}+\alpha$. For the
normalized equations also set $s=Y/q^3$ and
$n=n_{\rm old}+s$. These are positive. The exact substitutions give
\[
 Z_{\rm old}=q^3Z_{\rm new},\quad
 M_{\rm old}=q^3M_{\rm new},\quad
 N_{\rm old}=q^6N_{\rm new},
\]
so the new polynomial vanishes. Soundness was proved directly because
the inverse coordinate shift need not preserve positivity.

The three live operation counts are respectively
$(50,18,15)$, $(49,18,14)$, and $(48,18,14)$ for multiplication,
addition, and subtraction. Put $q_0=\beta J$ and $k_0=\eta+\zeta$.
Their highest homogeneous forms are
\begin{align*}
 P_{83}^{\rm top}&=2^{14}q_0^{200}w^{46}Y^{104}i^{20}k_0^{60}
                  (q_0-F)^4(q_0-F-\alpha)^2,\\
 P_{81}^{\rm top}&=-2^{10}q_0^{384}w^{34}s^{80}i^{14}k_0^{46}T^2,\\
 P_{80}^{\rm top}&=2^{14}q_0^{506}w^{46}s^{108}i^{20}k_0^{64}T^4
                  (q_0-F-\alpha)^2.
\end{align*}
They are nonzero for the authentic tuple and have degrees 436, 560,
and 750. Global gate-by-gate upper bounds and full univariate polynomial
replays independently verify attainment. A finite ray calculation alone
would not supply the global upper bound; both are included.


\section{Two compact unrestricted-integer equations}\label{sec:signed}
This component retains the entire six-integer source34 tuple of package 132,
including $\beta=2^{1064921}-1$, $C_A=2^{1113}$, and $C_B=2^{666104}$.
The other three numerals $K_0,M_C,M_F$ and the full finite reconstruction
recipe are supplied in \filepath{integer11/coefficients/}. These constants
are independent of the four-coefficient tuple in Section~\ref{sec:compact}.

The eleven unknowns $j,C,z,h,r,v,w,T,c,n,s$ range over all of $\Z$. Set
\begin{align*}
 J&=6j+1,&q&=\beta J+1,&W&=12w+4,\\
 X&=qW,&Y&=(2v+1)q^3,&E&=XY,\\
 a&=E+Y,&A&=a+2,&H&=4a+3,\\
 \Delta&=a^2+H=A^2-1,\\
 F&=(K_0+W)C-z(q-1),\\
 R&=(q(q-F)-C)(q^2-1)+(M_C+qM_F)J.
\end{align*}
Choose $M_0=E$ for the native version or $M_0=4q^4$ for the short-index
version, and continue with
\begin{align*}
 k&=R+1+hM_0,&\eta&=c-Yk,\\
 D_{\rm err}&=(C_AF-q)^2+(C_BC-q)^2,\\
 B_{\rm ud}&=k^2(q^2-D_{\rm err})-4q^2\eta^2,\\
 L&=XY^2k,&D_t&=1+L(L+k),&D&=B_{\rm ud}-4,\\
 P&=D_t+D(D_ts)^2,&D_m&=1+\Delta c^2,&B_P&=D_mP,\\
 U&=Arc^2,&D_f&=1+\Delta U^2B_P,&Q&=\Delta^2U^2B_P,\\
 V&=D_f(c-R+4UB_PT)-c,&K_+&=X(3c-2X)+2,\\
 S&=srH,&z_c&=HT+K_+,\\
 A_c&=(z_c^2+3H^2)(z_c^3-19H^3),&Z&=Sn-A_c,\\
 R_G&=S^2B_PD_f.
\end{align*}
The two final equations have the same shape
\begin{equation}\label{eq:signed}
 \boxed{Q(Z^2-R_GV^2)+(R_G-Z^2)=0.}
\end{equation}

The change combines package 132's main-root saturation with package 131's
parity guard. The authenticated affine front gives even $q$ and makes
$U_c=C_AF$, $V_c=C_BC$ divisible by sixteen. Expanding,
\[
 B_{\rm ud}+(qk)^2
 =k^2\bigl[2q(U_c+V_c)-U_c^2-V_c^2\bigr]-4q^2\eta^2
 \equiv0\pmod{16}.
\]
Thus $B_{\rm ud}\equiv0$ or $12\pmod{16}$, and
$D\equiv12$ or $8\pmod{16}$. In particular $D$ is nonsquare and
divisible by four. Since $D_t\ge1$ and $D_t\equiv1\pmod8$,
\[
 D_g=1+DD_ts^2\equiv1\pmod4,\qquad P=D_tD_g\equiv1\pmod4.
\]
The old factors three and four in the guard are unnecessary.

For completeness of the sign-sensitive argument, put
$\delta=Q-1$, $C_y=QV^2-1$, and $D_a=\delta C_y$.
The affine front has $\Delta>1$, nonzero even $A$, and nonzero odd $H$.
The polynomial $(u^2+3)(u^3-19)$ has no rational zero; consequently
$A_c\ne0$, and a zero of \eqref{eq:signed} forces $S\ne0$.
The three integers $B_P,D_f,D_a$ are nonzero and pairwise coprime.
Indeed $D_f\equiv D_a\equiv1\pmod{B_P}$, while modulo $D_f$
\[
 Q\equiv-\Delta,\quad V\equiv-c,\quad
 D_a\equiv A^2D_m.
\]
Because $A\mid U$ and $D_m\mid B_P$, $D_f$ is coprime to $A^2D_m$.
Multiplying the final equation by $\delta$ gives
\[
 (\delta Z/S)^2=B_PD_fD_a.
\]
The absolute values of the factors are therefore squares. Since
$D_f\equiv1\pmod4$, it is positive. If $c\ne0$, a negative $B_P$ would
make $D_f<0$. If $c=0$, then $D_m=D_f=D_a=1$, so the same square identity
gives $B_P=P>0$ directly. Hence $B_P,P,D_a>0$ in all cases.

Now $D_g>0$. If $B_{\rm ud}\le3$, then $D\le-1$ and $s\ne0$ would
give $D_g\le1-D_t\le0$. Thus $B_{\rm ud}\ge4$, in fact at least twelve
by its residue class. The centered budget recovers the strict front signs
and excludes $c=0$. Extract $p_0^2=B_P$ and define $I=Up_0$,
$T_{\rm old}=p_0T$. This restores the strong and auxiliary Pell norms,
with $c^2\mid I$. The parent's early folding theorem uses these norms
before the first Pell root: it gives $D_m=d^2$. Only then infer
$P=(p_0/d)^2$ and split the coprime positive factors $D_t,D_g$ into
squares. This order avoids circular root extraction.

The remaining signed rank, projector, and fresh source reconstruction are
the unchanged package-132 interfaces. For the converse, canonical outer
data have $D_t=\tau^2$ and $B_{\rm ud}>64$. Then $DD_t>0$ is nonsquare,
so Pell's theorem supplies $s,g>0$ with $g^2-DD_ts^2=1$. Fix
$p_0=d\tau g$ first, choose one strong pair with
$4sAp_0c^2\mid I$, and use the parent's supported-modulus interpolation
with that same pair. This constructs the required integer $T$ and $n$.
The interpolation depends on divisibility support, not on the removed
scalar factors in $D$.

Replacing $12B_{\rm ud}-16$ by $B_{\rm ud}-4$ saves one multiplication.
The exact finalizer identity
\[
 (Q-1)Z^2-R_G(QV^2-1)
 =Q(Z^2-R_GV^2)+(R_G-Z^2)
\]
saves one further operation. The complete circuits consequently have 107
and 108 operations at unchanged exact degrees 924 and 840. Their global
leading monomials, including nonzero authentic coefficients, are proved
by the independent symbolic checker. The new guard and finalizer have
also received a separate mathematical review in the archive.


\section{Reproduction, limitations, and further research}\label{sec:reproduce}
\subsection{What is supplied and what was checked}

The archive contains this PDF and all its LaTeX inputs, the principal
integer polynomial and its 319-variable companion, both natural quartics,
the three compact positive circuits, and both eleven-integer circuits.
Every endpoint has its ordered variable list, literal coefficients or
complete arithmetic DAG, source binding, and an evaluator or exact
verification entry point. Imported source proofs and licenses retain
their provenance.

The top-level command
\begin{center}\vtt{python3 -B verify_all.py --rebuild}\end{center}
runs the independent component checks and deterministic rebuilds from the
archive alone, in ordinary and optimized Python. Python's standard library
is sufficient. No network, proof search, machine simulation, symbolic
algebra package, or original repository checkout is required. The PDF
build is separate and uses pdfLaTeX through the supplied build script.

The finite checks cover the following distinct obligations:
\begin{itemize}
\item The direct integer source changes are checked literally, including
the nine powers at base $q+1$, the remaining fourteen guards, and the two
new square-coordinate rows. The 217 monic elimination steps, translations,
and all 6,607 collected coefficients of the principal equation are replayed.
The parity companion has its own complete replay.
\item The natural branch checks every changed and unchanged raw row,
both simultaneous coordinate shifts, all 101 elimination certificates per
endpoint, and all quartic coefficients. Its isolated producer regenerates
the twelve normative JSON files exactly. Four corrupted artifacts per
endpoint are rejected.
\item The positive branch authenticates the entire fixed coefficient tuple,
checks ten exact free-polynomial identities, every graph reference and live
operation, global degree bounds, and complete univariate polynomial
specializations with nonzero leading coefficients. Twelve actual gate
mutations are rejected. The final component receipt records 276 checks.
\item The eleven-integer branch reconstructs all six fixed coefficients
from their finite recipe, checks local identities and every live gate,
proves the full leading monomial over the coefficient ring, and rejects
literal, input, and degree corruptions.
\end{itemize}

These checks are not numerical evidence for consistency. They certify
finite algebra, exact artifacts, and transformations. The unbounded
Pell, factorial, CRT, and proof-calculus arguments are the ordinary
mathematical proofs in this report and the identified inherited sources.
Separate agents reviewed the new sign-sensitive arguments; those reviews
are internal independent checks in this AI-assisted workflow, not outside
human endorsements. The work does not claim a newly compiled Lean/Rocq
development or a sealed object-language ZFC proof.

\subsection{Comparison and scope}

The repository was inspected at
\begin{center}\small
\vtt{1c5367452a0bd4dd7e9e4bceb1508b62dbd2ff61}.
\end{center}
It includes packages 127--133. This matters because an older summary table
still listed weaker endpoints. The comparisons in Section 1 use literal
equations and live circuits from the newest relevant packages, not the
stale table. A complete manuscript or equation from an earlier package is
credited wherever its proof interface is reused.

The direct 286-variable equation is the recommended unrestricted-integer
result of this iteration. It reduces dimension substantially while keeping
every coefficient below $2^{17}$ in absolute value. It is not the
fewest-variable quartic in the repository. The natural 235-variable form
has a different domain and much larger coefficients; it must not be
silently presented as an equation in 235 unrestricted integers. Likewise,
the 80-operation compact form has ten \emph{positive} inputs and larger
degree than the 83-operation form. The ordinary ten-integer constructions
using Sun's reduction remain a different dimension/degree/constant-size
tradeoff~\cite{sun}.

The source repository is unchanged by this work. The deliverable is an
integration-ready archive, with a suggested destination in the active
Diophantine consistency research collection and a source ledger that pins
all comparisons.

\subsection{Further research questions}

\begin{enumerate}
\item \textbf{Can other arithmetic guards be recovered semantically?}
The principal saving came from using a power certificate twice: once for
its value and once for its exponent's sign. A systematic dependency audit
of the remaining fourteen guards could find further removals. Each
candidate must be proved before the step that needs it; a circular
factorial or radix argument would introduce spurious roots.
\item \textbf{Can the natural one-row cap be transferred to the signed
compiler?} The natural bound uses a nonnegative quotient essentially.
A signed analogue would need to recover that quotient's sign or replace
the bound with another exact arithmetic condition. Merely reusing the
same equation over $\Z$ is unsound.
\item \textbf{Can abscissa orientation be made sparser?} The pooled
$(x_i-1)^2$ terms buy 27 deleted sign witnesses but account for the
principal support increase. A smaller subset of principal orientations,
followed by propagation along shared power relations, may retain most of
the dimensional saving with fewer collected monomials.
\item \textbf{Can the 80-operation construction retain a lower degree?}
The $q^3$ normalization removes three gates but raises the degrees of
the Pell data. A different quotient coordinate or a compatible lower
degree carrier would improve this tradeoff. The leading forms in
Section~\ref{sec:compact} identify the precise terms to target.
\item \textbf{Can the new arithmetic interface be formalized independently
of the whole compiler?} The negative-exponent lemma, even-frame anchor,
linear cofinal bound, and free auxiliary-coordinate transport are compact
targets for Lean or Rocq. Proving these first would make subsequent
compiler changes easier to audit. A full internal-ZFC derivation remains
a separate task.
\end{enumerate}

The most useful general lesson of this iteration is constructive: before
adding a sign witness, inspect whether an already required arithmetic
relation and its size bound force that sign. The new equations make this
principle precise in an encoding of finite set-theoretic refutations.


\begin{thebibliography}{9}
\raggedright
\bibitem{proveit}
Vladimir Reshetnikov and contributors, \emph{ProveIt}, inspected commit
\vtt{1c5367452a0bd4dd7e9e4bceb1508b62dbd2ff61}.
\url{https://github.com/VladimirReshetnikov/ProveIt}.
The archive's source ledger supplies exact paths and hashes.
\bibitem{p131}
ProveIt package 131, \emph{Shared Pell caps and coprime product equations
for ZFC consistency}, direct and scalar components, at the inspected commit.
\bibitem{p132}
ProveIt package 132, \emph{Signed Pell caps and free scale carriers},
eleven-integer main-root saturation component, at the inspected commit.
\bibitem{p133}
ProveIt package 133, \emph{Staged Pell certificates and a product-square
equation for $\Con(\ZFC)$}, direct and compact components, at the inspected commit.
\bibitem{p128}
ProveIt package 128, \emph{One-row Pell certificates}, natural cap construction,
at the inspected commit.
\bibitem{mr}
Yuri Matiyasevich and Julia Robinson, ``Reduction of an arbitrary
Diophantine equation to one in 13 unknowns,'' \emph{Acta Arithmetica}
27 (1975), 521--553. \url{https://doi.org/10.4064/aa-27-1-521-553}.
\bibitem{sun}
Zhi-Wei Sun, \emph{Ten unknowns for Hilbert's tenth problem over the integers},
arXiv:2609.01594 (2026). \url{https://arxiv.org/abs/2609.01594}.
This is background for the dimension frontier, not a dependency of the
new direct quartics.
\end{thebibliography}

\clearpage\appendix
\section{Every residual of the 286-integer polynomial}\label{app:residuals}
The independent unknowns are $z_0,\ldots,z_{285}\in\Z$. The following
literal table, together with \eqref{eq:main}, completely defines the
principal polynomial. A repeated variable in a term denotes its ordinary
power. The ordered coordinate names and compiler restoration maps are
additional explanatory data; they are not needed to interpret this table.
% Generated from SHA-256 5039d1907601db05946a6535a3f83edcb6e96c5b2469e8062fa99472ec32c6d1
\begingroup\small
\setlength{\jot}{2pt}
\setlength{\abovedisplayskip}{3pt}
\setlength{\belowdisplayskip}{3pt}
\begin{align*}
r_{0}={}&-z_{3}^{2} +z_{4}
\end{align*}
\begin{align*}
r_{1}={}&-1 -z_{0}z_{2} +z_{8}
\end{align*}
\begin{align*}
r_{2}={}&-1 +z_{1} +z_{2} -z_{2}z_{6} -z_{3}\\
&-z_{6}
\end{align*}
\begin{align*}
r_{3}={}&-2 +z_{1} -z_{7}z_{8}
\end{align*}
\begin{align*}
r_{4}={}&1 +z_{4} -z_{5}z_{19} +z_{17}
\end{align*}
\begin{align*}
r_{5}={}&-1 +2z_{9} -z_{9}^{2} -2z_{9}z_{10} +2z_{10}\\
&-z_{10}^{2} +z_{20}
\end{align*}
\begin{align*}
r_{6}={}&-5 -30z_{9} -36z_{9}^{2} -96z_{9}z_{10} -24z_{9}z_{20}\\
&-40z_{10} -64z_{10}^{2} -32z_{10}z_{20} -10z_{20} -4z_{20}^{2}\\
&+z_{21}
\end{align*}
\begin{align*}
r_{7}={}&-2 +2z_{9} -z_{9}^{2} -4z_{9}z_{21} +4z_{21}\\
&-4z_{21}^{2} +z_{22}
\end{align*}
\begin{align*}
r_{8}={}&1 +6z_{9} -36z_{9}^{2} -192z_{9}z_{21} -24z_{9}z_{22}\\
&+16z_{21} -256z_{21}^{2} -64z_{21}z_{22} +2z_{22} -4z_{22}^{2}\\
&+z_{23}
\end{align*}
\begin{align*}
r_{9}={}&-2 -z_{10}^{2} -2z_{10}z_{11} -z_{11}^{2} +z_{24}
\end{align*}
\begin{align*}
r_{10}={}&-9 +7z_{9} -z_{9}^{2} -4z_{9}z_{10} -4z_{9}z_{24}\\
&+12z_{10} -4z_{10}^{2} -8z_{10}z_{24} +12z_{24} -4z_{24}^{2}\\
&+z_{25}
\end{align*}
\begin{align*}
r_{11}={}&-1 +2z_{9} -z_{9}^{2} -2z_{9}z_{11} +2z_{11}\\
&-z_{11}^{2} +z_{26}
\end{align*}
\begin{align*}
r_{12}={}&-6z_{9} -144z_{9}^{2} -192z_{9}z_{10} -192z_{9}z_{11} -48z_{9}z_{20}\\
&-48z_{9}z_{26} -8z_{10} -64z_{10}^{2} -128z_{10}z_{11} -32z_{10}z_{20}\\
&-32z_{10}z_{26} -64z_{11}^{2} -32z_{11}z_{20} -32z_{11}z_{26} -2z_{20}\\
&-4z_{20}^{2} -8z_{20}z_{26} -4z_{26}^{2} +z_{27}
\end{align*}
\begin{align*}
r_{13}={}&-4z_{25}^{2} -8z_{25}z_{27} -4z_{27}^{2} +z_{28}
\end{align*}
\begin{align*}
r_{14}={}&-6 +16z_{11} -16z_{11}^{2} -16z_{11}z_{24} -16z_{11}z_{26}\\
&+8z_{24} -4z_{24}^{2} -8z_{24}z_{26} +8z_{26} -4z_{26}^{2}\\
&+z_{29}
\end{align*}
\begin{align*}
r_{15}={}&1 +6z_{9} -36z_{9}^{2} -96z_{9}z_{10} -48z_{9}z_{11}\\
&-24z_{9}z_{20} -48z_{9}z_{26} -24z_{9}z_{29} +8z_{10} -64z_{10}^{2}\\
&-64z_{10}z_{11} -32z_{10}z_{20} -64z_{10}z_{26} -32z_{10}z_{29} +8z_{11}\\
&-16z_{11}^{2} -16z_{11}z_{20} -32z_{11}z_{26} -16z_{11}z_{29} +2z_{20}\\
&-4z_{20}^{2} -16z_{20}z_{26} -8z_{20}z_{29} +8z_{26} -16z_{26}^{2}\\
&-16z_{26}z_{29} +4z_{29} -4z_{29}^{2} +z_{30}
\end{align*}
\begin{align*}
r_{16}={}&-10 -5z_{9} -z_{9}^{2} -12z_{9}z_{10} -16z_{9}z_{11}\\
&-4z_{9}z_{24} -36z_{10} -36z_{10}^{2} -96z_{10}z_{11} -24z_{10}z_{24}\\
&-48z_{11} -64z_{11}^{2} -32z_{11}z_{24} -12z_{24} -4z_{24}^{2}\\
&+z_{31}
\end{align*}
\begin{align*}
r_{17}={}&-2 -24z_{9} -36z_{9}^{2} -12z_{9}z_{10} -96z_{9}z_{11}\\
&-24z_{9}z_{26} -3z_{10} -z_{10}^{2} -16z_{10}z_{11} -4z_{10}z_{26}\\
&-32z_{11} -64z_{11}^{2} -32z_{11}z_{26} -8z_{26} -4z_{26}^{2}\\
&+z_{32}
\end{align*}
\begin{align*}
r_{18}={}&2 -4z_{31}^{2} -8z_{31}z_{32} -4z_{32}^{2} +z_{33}
\end{align*}
\begin{align*}
r_{19}={}&-12 +16z_{9} -4z_{9}^{2} -8z_{9}z_{11} -8z_{9}z_{20}\\
&-8z_{9}z_{24} +16z_{11} -4z_{11}^{2} -8z_{11}z_{20} -8z_{11}z_{24}\\
&+16z_{20} -4z_{20}^{2} -8z_{20}z_{24} +16z_{24} -4z_{24}^{2}\\
&+z_{34}
\end{align*}
\begin{align*}
r_{20}={}&-3 -24z_{9} -64z_{9}^{2} -160z_{9}z_{11} -64z_{9}z_{24}\\
&-32z_{9}z_{26} -32z_{9}z_{34} -26z_{11} -100z_{11}^{2} -80z_{11}z_{24}\\
&-40z_{11}z_{26} -40z_{11}z_{34} -8z_{24} -16z_{24}^{2} -16z_{24}z_{26}\\
&-16z_{24}z_{34} -6z_{26} -4z_{26}^{2} -8z_{26}z_{34} -4z_{34}\\
&-4z_{34}^{2} +z_{35}
\end{align*}
\begin{align*}
r_{21}={}&-1 -18z_{9} -36z_{9}^{2} -96z_{9}z_{10} -24z_{9}z_{20}\\
&-24z_{10} -64z_{10}^{2} -32z_{10}z_{20} -6z_{20} -4z_{20}^{2}\\
&+z_{36}
\end{align*}
\begin{align*}
r_{22}={}&-2 +2z_{9} -z_{9}^{2} -4z_{9}z_{36} +4z_{36}\\
&-4z_{36}^{2} +z_{37}
\end{align*}
\begin{align*}
r_{23}={}&1 +6z_{9} -36z_{9}^{2} -192z_{9}z_{36} -24z_{9}z_{37}\\
&+16z_{36} -256z_{36}^{2} -64z_{36}z_{37} +2z_{37} -4z_{37}^{2}\\
&+z_{38}
\end{align*}
\begin{align*}
r_{24}={}&-z_{10}^{2} -4z_{10}z_{38} -4z_{38}^{2} +z_{39}
\end{align*}
\begin{align*}
r_{25}={}&-6 -30z_{9} -64z_{9}^{2} -128z_{9}z_{10} -64z_{9}z_{20}\\
&-32z_{10} -64z_{10}^{2} -64z_{10}z_{20} -14z_{20} -16z_{20}^{2}\\
&+z_{40}
\end{align*}
\begin{align*}
r_{26}={}&-4 +8z_{10} -64z_{10}^{2} -96z_{10}z_{11} -32z_{10}z_{24}\\
&-32z_{10}z_{40} +6z_{11} -36z_{11}^{2} -24z_{11}z_{24} -24z_{11}z_{40}\\
&+2z_{24} -4z_{24}^{2} -8z_{24}z_{40} +4z_{40} -4z_{40}^{2}\\
&+z_{41}
\end{align*}
\begin{align*}
r_{27}={}&8z_{9} -64z_{9}^{2} -96z_{9}z_{11} -32z_{9}z_{26} -32z_{9}z_{41}\\
&+6z_{11} -36z_{11}^{2} -24z_{11}z_{26} -24z_{11}z_{41} +2z_{26}\\
&-4z_{26}^{2} -8z_{26}z_{41} +4z_{41} -4z_{41}^{2} +z_{42}
\end{align*}
\begin{align*}
r_{28}={}&-z_{11}^{2} -2z_{11}z_{12} -z_{12}^{2} +z_{43}
\end{align*}
\begin{align*}
r_{29}={}&2z_{9} -z_{9}^{2} -4z_{9}z_{11} -4z_{9}z_{43} +4z_{11}\\
&-4z_{11}^{2} -8z_{11}z_{43} +4z_{43} -4z_{43}^{2} +z_{44}
\end{align*}
\begin{align*}
r_{30}={}&-35 +24z_{9} -4z_{9}^{2} -8z_{9}z_{10} -8z_{9}z_{24}\\
&-8z_{9}z_{26} +24z_{10} -4z_{10}^{2} -8z_{10}z_{24} -8z_{10}z_{26}\\
&+24z_{24} -4z_{24}^{2} -8z_{24}z_{26} +24z_{26} -4z_{26}^{2}\\
&+z_{45}
\end{align*}
\begin{align*}
r_{31}={}&-16z_{9}^{2} -32z_{9}z_{11} -32z_{9}z_{26} -32z_{9}z_{43} -16z_{9}z_{44}\\
&-16z_{9}z_{45} -4z_{11} -16z_{11}^{2} -32z_{11}z_{26} -32z_{11}z_{43}\\
&-16z_{11}z_{44} -16z_{11}z_{45} -16z_{26}^{2} -32z_{26}z_{43} -16z_{26}z_{44}\\
&-16z_{26}z_{45} -4z_{43} -16z_{43}^{2} -16z_{43}z_{44} -16z_{43}z_{45}\\
&-2z_{44} -4z_{44}^{2} -8z_{44}z_{45} -4z_{45}^{2} +z_{46}
\end{align*}
\begin{align*}
r_{32}={}&z_{11} -25z_{11}^{2} -40z_{11}z_{26} -20z_{11}z_{29} -16z_{26}^{2}\\
&-16z_{26}z_{29} -4z_{29}^{2} +z_{47}
\end{align*}
\begin{align*}
r_{33}={}&-1 +8z_{10} -4z_{10}^{2} -8z_{10}z_{24} +8z_{24}\\
&-4z_{24}^{2} +z_{48}
\end{align*}
\begin{align*}
r_{34}={}&-1 -16z_{10}^{2} -8z_{10}z_{11} -32z_{10}z_{24} -16z_{10}z_{48}\\
&-z_{11}^{2} -8z_{11}z_{24} -4z_{11}z_{48} -16z_{24}^{2} -16z_{24}z_{48}\\
&-4z_{48}^{2} +z_{49}
\end{align*}
\begin{align*}
r_{35}={}&1 +z_{10} -z_{10}^{2} -4z_{10}z_{11} -4z_{10}z_{49}\\
&-4z_{11}^{2} -8z_{11}z_{49} -4z_{49}^{2} +z_{50}
\end{align*}
\begin{align*}
r_{36}={}&-1 -4z_{50} -4z_{50}^{2} +z_{51}
\end{align*}
\begin{align*}
r_{37}={}&-1 -z_{11}^{2} -8z_{11}z_{50} -4z_{11}z_{51} -16z_{50}^{2}\\
&-16z_{50}z_{51} -4z_{51}^{2} +z_{52}
\end{align*}
\begin{align*}
r_{38}={}&2 +z_{10} -z_{10}^{2} -4z_{10}z_{11} -4z_{10}z_{52}\\
&-4z_{11}^{2} -8z_{11}z_{52} -4z_{52}^{2} +z_{53}
\end{align*}
\begin{align*}
r_{39}={}&-4z_{47}^{2} -8z_{47}z_{53} -4z_{53}^{2} +z_{54}
\end{align*}
\begin{align*}
r_{40}={}&-9 +7z_{9} -z_{9}^{2} -8z_{9}z_{47} -4z_{9}z_{54}\\
&+24z_{47} -16z_{47}^{2} -16z_{47}z_{54} +12z_{54} -4z_{54}^{2}\\
&+z_{55}
\end{align*}
\begin{align*}
r_{41}={}&-1 -4z_{55} -4z_{55}^{2} +z_{56}
\end{align*}
\begin{align*}
r_{42}={}&-z_{11}^{2} -8z_{11}z_{55} -4z_{11}z_{56} -16z_{55}^{2} -16z_{55}z_{56}\\
&-4z_{56}^{2} +z_{57}
\end{align*}
\begin{align*}
r_{43}={}&-z_{9} -z_{9}^{2} -4z_{9}z_{11} -4z_{9}z_{57} -4z_{11}\\
&-4z_{11}^{2} -8z_{11}z_{57} -4z_{57} -4z_{57}^{2} +z_{58}
\end{align*}
\begin{align*}
r_{44}={}&-10 -12z_{58} -4z_{58}^{2} +z_{59}
\end{align*}
\begin{align*}
r_{45}={}&1 +4z_{10} -4z_{10}^{2} -8z_{10}z_{24} -16z_{10}z_{58}\\
&-8z_{10}z_{59} +4z_{24} -4z_{24}^{2} -16z_{24}z_{58} -8z_{24}z_{59}\\
&+8z_{58} -16z_{58}^{2} -16z_{58}z_{59} +4z_{59} -4z_{59}^{2}\\
&+z_{60}
\end{align*}
\begin{align*}
r_{46}={}&z_{10} -25z_{10}^{2} -40z_{10}z_{24} -20z_{10}z_{60} -16z_{24}^{2}\\
&-16z_{24}z_{60} -4z_{60}^{2} +z_{61}
\end{align*}
\begin{align*}
r_{47}={}&-10 -12z_{53} -4z_{53}^{2} +z_{62}
\end{align*}
\begin{align*}
r_{48}={}&1 -16z_{53}^{2} -16z_{53}z_{61} -16z_{53}z_{62} -4z_{61}^{2}\\
&-8z_{61}z_{62} -4z_{62}^{2} +z_{63}
\end{align*}
\begin{align*}
r_{49}={}&z_{11} -z_{11}^{2} -8z_{11}z_{61} -4z_{11}z_{63} -16z_{61}^{2}\\
&-16z_{61}z_{63} -4z_{63}^{2} +z_{64}
\end{align*}
\begin{align*}
r_{50}={}&-4z_{9} -4z_{9}^{2} -4z_{9}z_{11} -8z_{9}z_{20} -z_{11}\\
&-z_{11}^{2} -4z_{11}z_{20} -4z_{20} -4z_{20}^{2} +z_{65}
\end{align*}
\begin{align*}
r_{51}={}&1 +4z_{9} -4z_{9}^{2} -8z_{9}z_{20} -8z_{9}z_{65}\\
&+4z_{20} -4z_{20}^{2} -8z_{20}z_{65} +4z_{65} -4z_{65}^{2}\\
&+z_{66}
\end{align*}
\begin{align*}
r_{52}={}&-2 -4z_{9} -16z_{9}^{2} -32z_{9}z_{20} -16z_{9}z_{64}\\
&-16z_{9}z_{66} -4z_{20} -16z_{20}^{2} -16z_{20}z_{64} -16z_{20}z_{66}\\
&-4z_{64}^{2} -8z_{64}z_{66} -2z_{66} -4z_{66}^{2} +z_{67}
\end{align*}
\begin{align*}
r_{53}={}&-2 -4z_{67} -4z_{67}^{2} +z_{68}
\end{align*}
\begin{align*}
r_{54}={}&-z_{3}z_{14} +z_{69}
\end{align*}
\begin{align*}
r_{55}={}&z_{13} -z_{14} -2z_{23}z_{69} -z_{70}
\end{align*}
\begin{align*}
r_{56}={}&z_{13} -z_{14} -4z_{25}z_{69} -2z_{28}z_{69} -z_{71}
\end{align*}
\begin{align*}
r_{57}={}&z_{13} -z_{14} -2z_{30}z_{69} -6z_{69} -z_{72}
\end{align*}
\begin{align*}
r_{58}={}&z_{13} -z_{14} -4z_{31}z_{69} -2z_{33}z_{69} -z_{73}
\end{align*}
\begin{align*}
r_{59}={}&z_{13} -z_{14} -2z_{35}z_{69} +6z_{69} -z_{74}
\end{align*}
\begin{align*}
r_{60}={}&z_{13} -z_{14} -4z_{38}z_{69} -2z_{39}z_{69} -z_{75}
\end{align*}
\begin{align*}
r_{61}={}&z_{13} -z_{14} -2z_{42}z_{69} -z_{76}
\end{align*}
\begin{align*}
r_{62}={}&z_{13} -z_{14} -2z_{46}z_{69} -4z_{69} -z_{77}
\end{align*}
\begin{align*}
r_{63}={}&-z_{9}z_{10} +z_{9}z_{11} +z_{10} +z_{10}z_{11} -z_{11}\\
&-z_{11}^{2} +z_{78}
\end{align*}
\begin{align*}
r_{64}={}&1 +2z_{12} +z_{12}^{2} -z_{78}^{2} +z_{79}
\end{align*}
\begin{align*}
r_{65}={}&-z_{79}^{2} +z_{80}
\end{align*}
\begin{align*}
r_{66}={}&-4z_{67}z_{69} -2z_{68}z_{69} +2z_{69} +z_{81}
\end{align*}
\begin{align*}
r_{67}={}&-z_{13}^{2} +2z_{13}z_{14} +2z_{13}z_{81} -z_{14}^{2} -2z_{14}z_{81}\\
&-z_{81}^{2} +z_{82}
\end{align*}
\begin{align*}
r_{68}={}&-1 -z_{3}z_{10} +z_{83}
\end{align*}
\begin{align*}
r_{69}={}&-z_{14}z_{83} +z_{85}z_{86} +z_{87}
\end{align*}
\begin{align*}
r_{70}={}&-1 +z_{3} -z_{3}z_{9} +z_{84}
\end{align*}
\begin{align*}
r_{71}={}&-1 +2z_{3} -2z_{3}z_{9} -2z_{3}z_{20} +z_{85}
\end{align*}
\begin{align*}
r_{72}={}&-z_{13}z_{84} +z_{86}
\end{align*}
\begin{align*}
r_{73}={}&-1 -6z_{3}z_{9} -8z_{3}z_{10} -2z_{3}z_{20} +z_{88}
\end{align*}
\begin{align*}
r_{74}={}&z_{13}z_{83} -z_{14}z_{88} +z_{89}
\end{align*}
\begin{align*}
r_{75}={}&-z_{13}z_{70} +z_{14}z_{70} +z_{90}
\end{align*}
\begin{align*}
r_{76}={}&-z_{71}z_{90} +z_{91}
\end{align*}
\begin{align*}
r_{77}={}&-z_{72}z_{91} +z_{92}
\end{align*}
\begin{align*}
r_{78}={}&-z_{73}z_{92} +z_{93}
\end{align*}
\begin{align*}
r_{79}={}&-z_{74}z_{93} +z_{94}
\end{align*}
\begin{align*}
r_{80}={}&-z_{75}z_{94} +z_{95}
\end{align*}
\begin{align*}
r_{81}={}&-z_{76}z_{95} +z_{96}
\end{align*}
\begin{align*}
r_{82}={}&-z_{77}z_{96} +z_{97}
\end{align*}
\begin{align*}
r_{83}={}&-z_{87}z_{97} +z_{98}
\end{align*}
\begin{align*}
r_{84}={}&-z_{89}z_{98} +z_{99}
\end{align*}
\begin{align*}
r_{85}={}&-z_{99}^{2} +z_{100}
\end{align*}
\begin{align*}
r_{86}={}&-1 -z_{2} -z_{2}z_{17} +z_{101}
\end{align*}
\begin{align*}
r_{87}={}&-1 -z_{15}z_{101} +z_{102}
\end{align*}
\begin{align*}
r_{88}={}&-2z_{1}^{2} +2z_{1}z_{13} +2z_{1}z_{14} +2z_{1}z_{102} +2z_{1}z_{103}\\
&-z_{13}^{2} -2z_{13}z_{102} -z_{14}^{2} -2z_{14}z_{103} +z_{18}z_{105}\\
&-z_{80}z_{100} -z_{82}z_{100} -z_{102}^{2} -z_{103}^{2}
\end{align*}
\begin{align*}
r_{89}={}&-1 -z_{2}z_{16} -z_{16}z_{101} +z_{103}
\end{align*}
\begin{align*}
r_{90}={}&-9z_{9} -9z_{10} -4z_{20} -2z_{23} -4z_{25}\\
&-2z_{28} -2z_{30} -4z_{31} -2z_{33} -2z_{35}\\
&-4z_{38} -2z_{39} -2z_{42} -2z_{46} -4z_{67}\\
&-2z_{68} +z_{104}
\end{align*}
\begin{align*}
r_{91}={}&-z_{106}z_{109} +z_{109} +z_{111}
\end{align*}
\begin{align*}
r_{92}={}&-z_{5}z_{108} +z_{107} -z_{108}z_{111} -z_{110}
\end{align*}
\begin{align*}
r_{93}={}&-z_{18}z_{113} +z_{116}
\end{align*}
\begin{align*}
r_{94}={}&z_{112} -z_{114}z_{116} -z_{115}
\end{align*}
\begin{align*}
r_{95}={}&-z_{116}z_{118} +z_{117} -z_{119}
\end{align*}
\begin{align*}
r_{96}={}&-z_{116}z_{121} +z_{120} -z_{122}
\end{align*}
\begin{align*}
r_{97}={}&-z_{116}z_{124} +z_{123} -z_{125}
\end{align*}
\begin{align*}
r_{98}={}&-z_{18}z_{127} +z_{130}
\end{align*}
\begin{align*}
r_{99}={}&z_{126} -z_{128}z_{130} -z_{129}
\end{align*}
\begin{align*}
r_{100}={}&-z_{130}z_{132} +z_{131} -z_{133}
\end{align*}
\begin{align*}
r_{101}={}&-z_{18}z_{135} +z_{138}
\end{align*}
\begin{align*}
r_{102}={}&z_{134} -z_{136}z_{138} -z_{137}
\end{align*}
\begin{align*}
r_{103}={}&-z_{138}z_{140} +z_{139} -z_{141}
\end{align*}
\begin{align*}
r_{104}={}&-z_{18}z_{108} +z_{145}
\end{align*}
\begin{align*}
r_{105}={}&z_{142} -z_{143}z_{145} -z_{144}
\end{align*}
\begin{align*}
r_{106}={}&-z_{116}z_{147} +z_{146} -z_{148}
\end{align*}
\begin{align*}
r_{107}={}&z_{149} -2z_{149}z_{151} +z_{152}
\end{align*}
\begin{align*}
r_{108}={}&-z_{149}z_{152} -z_{152} +z_{153}
\end{align*}
\begin{align*}
r_{109}={}&2z_{150} +z_{150}^{2} +z_{152}^{2} -z_{153}^{2}
\end{align*}
\begin{align*}
r_{110}={}&1 +z_{150} -z_{150}z_{155} -z_{155} +z_{156}
\end{align*}
\begin{align*}
r_{111}={}&1 +2z_{154} +z_{154}^{2} -2z_{155} +z_{155}^{2}\\
&-z_{156}^{2}
\end{align*}
\begin{align*}
r_{112}={}&z_{4}^{2} -2z_{4}z_{150} -2z_{150} +z_{157}
\end{align*}
\begin{align*}
r_{113}={}&1 +z_{2} +z_{4} -z_{4}z_{155} -z_{154}\\
&-z_{155} +z_{156} +z_{157}z_{158}
\end{align*}
\begin{align*}
r_{114}={}&-z_{150}z_{160} -z_{160} +z_{161}
\end{align*}
\begin{align*}
r_{115}={}&2z_{159} +z_{159}^{2} +z_{160}^{2} -z_{161}^{2}
\end{align*}
\begin{align*}
r_{116}={}&1 -2z_{5} +z_{5}^{2} -2z_{5}z_{150} +z_{162}
\end{align*}
\begin{align*}
r_{117}={}&-1 -z_{5}z_{160} +z_{18} -z_{159} +z_{161}\\
&+z_{162}z_{163}
\end{align*}
\begin{align*}
r_{118}={}&-1 +2z_{150}z_{223} -z_{164} +2z_{223} -z_{223}^{2}
\end{align*}
\begin{align*}
r_{119}={}&-2 +z_{106} -z_{159} -z_{160}z_{223} +z_{161}\\
&+z_{164}z_{165}
\end{align*}
\begin{align*}
r_{120}={}&1 +z_{150} -z_{150}z_{167} -z_{167} +z_{168}
\end{align*}
\begin{align*}
r_{121}={}&1 +2z_{166} +z_{166}^{2} -2z_{167} +z_{167}^{2}\\
&-z_{168}^{2}
\end{align*}
\begin{align*}
r_{122}={}&1 -2z_{106} +z_{106}^{2} -2z_{106}z_{150} +z_{169}
\end{align*}
\begin{align*}
r_{123}={}&-1 +z_{106} -z_{106}z_{167} +z_{107} -z_{166}\\
&+z_{168} +z_{169}z_{170}
\end{align*}
\begin{align*}
r_{124}={}&4 -4z_{106} +z_{106}^{2} -2z_{106}z_{150} +2z_{150}\\
&+z_{171}
\end{align*}
\begin{align*}
r_{125}={}&-2 +z_{106} -z_{106}z_{155} +z_{108} -z_{154}\\
&+z_{155} +z_{156} +z_{171}z_{172}
\end{align*}
\begin{align*}
r_{126}={}&1 +z_{150} -z_{150}z_{174} -z_{174} +z_{175}
\end{align*}
\begin{align*}
r_{127}={}&1 +2z_{173} +z_{173}^{2} -2z_{174} +z_{174}^{2}\\
&-z_{175}^{2}
\end{align*}
\begin{align*}
r_{128}={}&-1 +z_{106} -z_{106}z_{174} +z_{112} +z_{169}z_{176}\\
&-z_{173} +z_{175}
\end{align*}
\begin{align*}
r_{129}={}&-1 -2z_{150} -z_{150}^{2} +z_{177}
\end{align*}
\begin{align*}
r_{130}={}&-z_{106} -z_{106}z_{154} -z_{106}z_{156} +z_{113} +z_{150}z_{156}\\
&+z_{154} +z_{155} -z_{155}z_{177} +2z_{156} +z_{171}z_{178}\\
&+z_{177}
\end{align*}
\begin{align*}
r_{131}={}&-z_{150}z_{180} -z_{180} +z_{181}
\end{align*}
\begin{align*}
r_{132}={}&2z_{179} +z_{179}^{2} +z_{180}^{2} -z_{181}^{2}
\end{align*}
\begin{align*}
r_{133}={}&-1 -z_{106}z_{180} +z_{117} +z_{169}z_{182} -z_{179}\\
&+z_{181}
\end{align*}
\begin{align*}
r_{134}={}&-z_{150}z_{184} -z_{184} +z_{185}
\end{align*}
\begin{align*}
r_{135}={}&2z_{183} +z_{183}^{2} +z_{184}^{2} -z_{185}^{2}
\end{align*}
\begin{align*}
r_{136}={}&-1 -z_{106}z_{184} +z_{120} +z_{169}z_{186} -z_{183}\\
&+z_{185}
\end{align*}
\begin{align*}
r_{137}={}&-z_{150}z_{188} -z_{188} +z_{189}
\end{align*}
\begin{align*}
r_{138}={}&2z_{187} +z_{187}^{2} +z_{188}^{2} -z_{189}^{2}
\end{align*}
\begin{align*}
r_{139}={}&-1 -z_{106}z_{188} +z_{123} +z_{169}z_{190} -z_{187}\\
&+z_{189}
\end{align*}
\begin{align*}
r_{140}={}&-z_{150}z_{192} -z_{192} +z_{193}
\end{align*}
\begin{align*}
r_{141}={}&2z_{191} +z_{191}^{2} +z_{192}^{2} -z_{193}^{2}
\end{align*}
\begin{align*}
r_{142}={}&-1 -z_{106}z_{192} +z_{126} +z_{169}z_{194} -z_{191}\\
&+z_{193}
\end{align*}
\begin{align*}
r_{143}={}&-1 -z_{150} -z_{150}z_{196} -z_{196} +z_{197}
\end{align*}
\begin{align*}
r_{144}={}&1 +2z_{195} +z_{195}^{2} +2z_{196} +z_{196}^{2}\\
&-z_{197}^{2}
\end{align*}
\begin{align*}
r_{145}={}&-z_{106} -z_{106}z_{196} +z_{127} +z_{171}z_{198} -z_{195}\\
&+z_{196} +z_{197}
\end{align*}
\begin{align*}
r_{146}={}&-z_{150}z_{200} -z_{200} +z_{201}
\end{align*}
\begin{align*}
r_{147}={}&2z_{199} +z_{199}^{2} +z_{200}^{2} -z_{201}^{2}
\end{align*}
\begin{align*}
r_{148}={}&-1 -z_{106}z_{200} +z_{131} +z_{169}z_{202} -z_{199}\\
&+z_{201}
\end{align*}
\begin{align*}
r_{149}={}&-z_{150}z_{204} -z_{204} +z_{205}
\end{align*}
\begin{align*}
r_{150}={}&2z_{203} +z_{203}^{2} +z_{204}^{2} -z_{205}^{2}
\end{align*}
\begin{align*}
r_{151}={}&-1 -z_{106}z_{204} +z_{134} +z_{169}z_{206} -z_{203}\\
&+z_{205}
\end{align*}
\begin{align*}
r_{152}={}&-z_{150}z_{208} -z_{208} +z_{209}
\end{align*}
\begin{align*}
r_{153}={}&2z_{207} +z_{207}^{2} +z_{208}^{2} -z_{209}^{2}
\end{align*}
\begin{align*}
r_{154}={}&-1 -z_{106}z_{208} +z_{135} +z_{171}z_{210} -z_{207}\\
&+z_{208} +z_{209}
\end{align*}
\begin{align*}
r_{155}={}&-z_{150}z_{212} -z_{212} +z_{213}
\end{align*}
\begin{align*}
r_{156}={}&2z_{211} +z_{211}^{2} +z_{212}^{2} -z_{213}^{2}
\end{align*}
\begin{align*}
r_{157}={}&-1 -z_{106}z_{212} +z_{139} +z_{169}z_{214} -z_{211}\\
&+z_{213}
\end{align*}
\begin{align*}
r_{158}={}&-z_{150}z_{216} -z_{216} +z_{217}
\end{align*}
\begin{align*}
r_{159}={}&2z_{215} +z_{215}^{2} +z_{216}^{2} -z_{217}^{2}
\end{align*}
\begin{align*}
r_{160}={}&-1 -z_{106}z_{216} +z_{142} +z_{169}z_{218} -z_{215}\\
&+z_{217}
\end{align*}
\begin{align*}
r_{161}={}&-z_{150}z_{220} -z_{220} +z_{221}
\end{align*}
\begin{align*}
r_{162}={}&2z_{219} +z_{219}^{2} +z_{220}^{2} -z_{221}^{2}
\end{align*}
\begin{align*}
r_{163}={}&-1 -z_{106}z_{220} +z_{146} +z_{169}z_{222} -z_{219}\\
&+z_{221}
\end{align*}
\begin{align*}
r_{164}={}&-16 +8z_{0} -14z_{0}^{2} -16z_{0}z_{1} -16z_{0}z_{2}\\
&-16z_{0}z_{3} -4z_{0}z_{4} -10z_{1}^{2} -16z_{1}z_{2} -16z_{1}z_{3}\\
&-8z_{2} -12z_{2}^{2} -16z_{2}z_{3} -8z_{3}^{2} -4z_{4}^{2}\\
&-2z_{5}^{2} +4z_{9} -2z_{9}^{2} -2z_{10}^{2} -2z_{11}^{2}\\
&-2z_{12}^{2} -2z_{13}^{2} -2z_{14}^{2} -2z_{15}^{2} -2z_{16}^{2}\\
&-2z_{17}^{2} -2z_{18}^{2} -2z_{104}^{2} +4z_{106} -4z_{106}^{2}\\
&-2z_{107}^{2} -2z_{108}^{2} -2z_{112}^{2} -2z_{113}^{2} -2z_{117}^{2}\\
&-2z_{120}^{2} -2z_{123}^{2} -2z_{126}^{2} -2z_{127}^{2} -2z_{131}^{2}\\
&-2z_{134}^{2} -2z_{135}^{2} -2z_{139}^{2} -2z_{142}^{2} -2z_{146}^{2}\\
&+z_{149} -2z_{223}^{2}
\end{align*}
\begin{align*}
r_{165}={}&-z_{5} -2z_{5}z_{9} -2z_{5}z_{10} -2z_{5}z_{11} -2z_{5}z_{12}\\
&-2z_{5}z_{13} -2z_{5}z_{14} -2z_{5}z_{15} -2z_{5}z_{16} -2z_{5}z_{17}\\
&+z_{223}
\end{align*}
\begin{align*}
r_{166}={}&-17 +z_{0} -z_{224}^{2}
\end{align*}
\begin{align*}
r_{167}={}&-z_{150}z_{226} -z_{226} +z_{237}
\end{align*}
\begin{align*}
r_{168}={}&-1 +z_{225}^{2} +z_{226}^{2} -z_{237}^{2}
\end{align*}
\begin{align*}
r_{169}={}&-z_{150}z_{228} -z_{228} +z_{238}
\end{align*}
\begin{align*}
r_{170}={}&-1 +z_{227}^{2} +z_{228}^{2} -z_{238}^{2}
\end{align*}
\begin{align*}
r_{171}={}&-z_{230} -z_{230}z_{231} +z_{239}
\end{align*}
\begin{align*}
r_{172}={}&-1 +z_{229}^{2} +z_{230}^{2} -z_{239}^{2}
\end{align*}
\begin{align*}
r_{173}={}&-z_{226}^{2} +z_{240}
\end{align*}
\begin{align*}
r_{174}={}&z_{228} -4z_{233}z_{240} -2z_{240}
\end{align*}
\begin{align*}
r_{175}={}&-4z_{228} -8z_{228}z_{232} +z_{231}
\end{align*}
\begin{align*}
r_{176}={}&-z_{150} -z_{227}z_{234} +z_{231}
\end{align*}
\begin{align*}
r_{177}={}&-z_{225} -z_{227}z_{235} +z_{229}
\end{align*}
\begin{align*}
r_{178}={}&-z_{149} -4z_{226}z_{236} +z_{230}
\end{align*}
\begin{align*}
r_{179}={}&6z_{0}^{2} +8z_{0}z_{1} +8z_{0}z_{2} +8z_{0}z_{3} +2z_{0}z_{4}\\
&-2z_{0}z_{155} -4z_{0}z_{160} -2z_{0}z_{167} +5z_{1}^{2} +8z_{1}z_{2}\\
&+8z_{1}z_{3} -4z_{1}z_{160} -2z_{1}z_{208} +5z_{2}^{2} +8z_{2}z_{3}\\
&-4z_{2}z_{160} -2z_{2}z_{196} +4z_{3}^{2} -4z_{3}z_{160} +z_{4}^{2}\\
&-2z_{4}z_{167} +z_{9}^{2} -2z_{9}z_{174} +z_{10}^{2} -2z_{10}z_{180}\\
&+z_{11}^{2} -2z_{11}z_{184} +z_{12}^{2} -2z_{12}z_{188} +z_{13}^{2}\\
&-2z_{13}z_{192} +z_{14}^{2} -2z_{14}z_{200} +z_{15}^{2} -2z_{15}z_{204}\\
&+z_{16}^{2} -2z_{16}z_{212} +z_{17}^{2} -2z_{17}z_{216} +z_{104}^{2}\\
&-2z_{104}z_{220} -z_{150}z_{241} +z_{154}^{2} +z_{155}^{2} +z_{159}^{2}\\
&+z_{160}^{2} +z_{166}^{2} +z_{167}^{2} +z_{173}^{2} +z_{174}^{2}\\
&+z_{179}^{2} +z_{180}^{2} +z_{183}^{2} +z_{184}^{2} +z_{187}^{2}\\
&+z_{188}^{2} +z_{191}^{2} +z_{192}^{2} +z_{195}^{2} +z_{196}^{2}\\
&+z_{199}^{2} +z_{200}^{2} +z_{203}^{2} +z_{204}^{2} +z_{207}^{2}\\
&+z_{208}^{2} +z_{211}^{2} +z_{212}^{2} +z_{215}^{2} +z_{216}^{2}\\
&+z_{219}^{2} +z_{220}^{2}
\end{align*}
\begin{align*}
r_{180}={}&2 +z_{150} -2z_{242}^{2}
\end{align*}
\begin{align*}
r_{181}={}&-2 +z_{106} -4z_{243} -4z_{243}^{2}
\end{align*}
\begin{align*}
r_{182}={}&z_{2} -z_{3}z_{5} +z_{5} -z_{244}^{2} -z_{245}^{2}\\
&-z_{246} -z_{246}^{2}
\end{align*}
\begin{align*}
r_{183}={}&-1 -z_{2} +z_{3}z_{5} -z_{247}^{2} -z_{248}^{2}\\
&-z_{249} -z_{249}^{2}
\end{align*}
\begin{align*}
r_{184}={}&z_{108}z_{110} -z_{110} -z_{110}^{2} -z_{250}^{2} -z_{251}^{2}\\
&-z_{252} -z_{252}^{2}
\end{align*}
\begin{align*}
r_{185}={}&z_{113}z_{115} -z_{115} -z_{115}^{2} -z_{253}^{2} -z_{254}^{2}\\
&-z_{255} -z_{255}^{2}
\end{align*}
\begin{align*}
r_{186}={}&z_{113}z_{119} -z_{119} -z_{119}^{2} -z_{256}^{2} -z_{257}^{2}\\
&-z_{258} -z_{258}^{2}
\end{align*}
\begin{align*}
r_{187}={}&z_{113}z_{122} -z_{122} -z_{122}^{2} -z_{259}^{2} -z_{260}^{2}\\
&-z_{261} -z_{261}^{2}
\end{align*}
\begin{align*}
r_{188}={}&z_{113}z_{125} -z_{125} -z_{125}^{2} -z_{262}^{2} -z_{263}^{2}\\
&-z_{264} -z_{264}^{2}
\end{align*}
\begin{align*}
r_{189}={}&z_{127}z_{129} -z_{129} -z_{129}^{2} -z_{265}^{2} -z_{266}^{2}\\
&-z_{267} -z_{267}^{2}
\end{align*}
\begin{align*}
r_{190}={}&z_{127}z_{133} -z_{133} -z_{133}^{2} -z_{268}^{2} -z_{269}^{2}\\
&-z_{270} -z_{270}^{2}
\end{align*}
\begin{align*}
r_{191}={}&z_{135}z_{137} -z_{137} -z_{137}^{2} -z_{271}^{2} -z_{272}^{2}\\
&-z_{273} -z_{273}^{2}
\end{align*}
\begin{align*}
r_{192}={}&z_{135}z_{141} -z_{141} -z_{141}^{2} -z_{274}^{2} -z_{275}^{2}\\
&-z_{276} -z_{276}^{2}
\end{align*}
\begin{align*}
r_{193}={}&z_{108}z_{144} -z_{144} -z_{144}^{2} -z_{277}^{2} -z_{278}^{2}\\
&-z_{279} -z_{279}^{2}
\end{align*}
\begin{align*}
r_{194}={}&z_{113}z_{148} -z_{148} -z_{148}^{2} -z_{280}^{2} -z_{281}^{2}\\
&-z_{282} -z_{282}^{2}
\end{align*}
\begin{align*}
r_{195}={}&-5 -z_{149} +2z_{155} -z_{155}^{2} -z_{160}^{2}\\
&+2z_{167} -z_{167}^{2} +2z_{174} -z_{174}^{2} -z_{180}^{2}\\
&-z_{184}^{2} -z_{188}^{2} -z_{192}^{2} -2z_{196} -z_{196}^{2}\\
&-z_{200}^{2} -z_{204}^{2} -z_{208}^{2} -z_{212}^{2} -z_{216}^{2}\\
&-z_{220}^{2} +z_{226} -z_{283}^{2} -z_{284}^{2} -z_{285}\\
&-z_{285}^{2}
\end{align*}
\endgroup


\section{Fixed constants of the positive compact equations}\label{app:constants}
The following is the parent's complete finite constant recipe, retained
with its fixed indexed transition table. In this appendix the symbol $K$
means the constant $K_0$ in Section~\ref{sec:compact}; the local constant
$B$ is a numeral base, not the polynomial transport expression named $B$.
All indices and summation limits here are fixed. Floors and rational
intermediates define constants only and are not operations of the
Diophantine equation. The exact integer payloads and hashes are supplied.

\providecommand{\fl}[1]{\left\lfloor#1\right\rfloor}
\providecommand{\rg}[2]{\{#1,\ldots,#2\}}
\textbf{The four fixed integers.}
Every index, set, sum, product and floor below defines a constant, not
another unknown. Put
\[
\rho=2^{67},\quad B=\rho^{8807},\quad\beta=B-1,\quad
G_m(x)=\sum_{j=0}^{m-1}x^j.
\]
Define
\[
\begin{split}
\mathcal E={}&\rg0{942}\cup\rg{983}{992}\cup\rg{1003}{1011}\\
&\cup\rg{1013}{1031}\cup\{1063,1064,1095,2195,3296\}.
\end{split}
\]
Then
\begin{align*}
M_C={}&B-1-\sum_{e\in\mathcal E\setminus\{1\}}\rho^e-2\rho^{1063},\\
M_F={}&B-1+4+(2^{51}-2)(\rho^{4393}+\rho^{5495})\\
&+\rho^{1013}G_{10}(\rho)+\rho^{1065}G_9(\rho)
  +\rho^{2195}+\rho^{3296}\\
&+(2^{65}-1)\rho^{8792}G_{13}(\rho)+(2^{19}-1)\rho^{8805},\\
K={}&\rho^{30}+\rho^{62}+\rho^{1100}+\rho^{2201}+\rho^{8849}\\
&+(\rho^{4393}+\rho^{5495})\Phi(1/\rho),
\end{align*}
where
\[
\begin{split}
\Phi(x)={}&(1+2^{51})G_{943}(x)
 -2^{10}\sum_{i=0}^{942}v(t_i)x^i\\
&+\sum_{(b,e,m)\in\mathcal C}2^b x^e G_m(2x),\\
\mathcal C={}&\{(10,1023,8),(18,1013,9),(27,1003,8),\\
&(35,983,9),(44,1031,1),(45,1022,1),(46,1011,1),\\
&(47,992,1),(48,1095,1),(49,2195,1),(50,3296,1)\}.
\end{split}
\]
The rational intermediate in K cancels exactly; all four results are positive integers.

\textbf{The 943 indexed fixed tiles.}
Use the actions $u_j,v_j,d_j$ decoded below. Let
\[
\begin{gathered}
h(q,b)=2(q+1)+b,\quad
\mathcal R=\{u_j:d_j=1,\ 0\le j<430\},\\
\mathcal L=\{u_j:d_j=0,\ 0\le j<430\},\\
r_q=6+|\mathcal R\cap\rg0{q-1}|,\quad
l_q=205+|\mathcal L\cap\rg0{q-1}|.
\end{gathered}
\]
An interval ending at $-1$ is empty. The table defines $t_i$ for
$0\le i\le942$; all index ranges are fixed:
\[
\setlength{\arraycolsep}{1.4pt}\begin{array}{c|l}
\text{index }i&\text{tile }t_i=(w,n,e,s)\\\hline
0,1&(2,435,3,1),\ (5,0,0,0)\\
j+2\ (j=2q+b<430)&((1-d_j)l_{u_j},h(q,b),d_jr_{u_j},v_j)\\
432+b\ (b=0,1)&(0,h(215,b),0,h(215,b))\\
434+b\ (b=0,1)&(0,b,0,b)\\
424+2r_q+b\ (q\in\mathcal R,b=0,1)&(r_q,b,0,h(q,b))\\
424+2l_q+b\ (q\in\mathcal L,b=0,1)&(0,b,l_q,h(q,b))\\
924+j\ (0\le j<4)&(0,j+430\fl{j/2},0,438)\\
928&(2,434,2,0)\\
929&(3,436,4,4)\\
930&(4,436,4,0)\\
497+t\ (t=434,435,436)&(1,438,1,t)\\
\substack{934+3a+b-6\fl{a/4}\\a\in\{0,1,4\},\ b=0,1,2}&(a,437,b,437)
\end{array}
\]
For $t=(w,n,e,s)$ put
\[
v(t)=w+2^8n+2^{17}e+2^{25}s+127\,2^{34}.
\]

\textbf{A closed decoder for the 430 actions.}
Put $D_b(A,k)=\fl{A/b^k}\bmod b$, where
$a\bmod b=a-b\fl{a/b}$. For $0\le q<215$, $b=0,1$, $j=2q+b$:
\[
\begin{gathered}
N=\mathcal A\bmod2^{636},\quad
T_0=\fl{\mathcal A/2^{636}}\bmod2^{1675},\\
E_0=\fl{\mathcal A/2^{2311}},\\
\nu_j=D_2(N,j),\quad \sigma_j=\sum_{k=0}^{j-1}\nu_k,\\
m_q=D_2(N,430+q),\quad \tau_q=\sum_{k=0}^{q-1}m_k,\\
F_q=14(1-m_q)+m_qD_{16}(E_0,\tau_q),\\
u_j=\nu_j(1+\sigma_j)+(1-\nu_j)D_{216}(T_0,j-\sigma_j),\\
v_j=D_2(F_q,2b),\quad d_j=D_2(F_q,2b+1).
\end{gathered}
\]
Empty sums are zero. Here $(v_0,d_0,u_0)=(1,1,1)$; the constructor tile
uses this actual first target, including its state renaming.

\textbf{The integer $\mathcal A$.}
The alphabet, in order of digit values $0,\ldots,61$, is\par\noindent
{\fontsize{9}{10}\selectfont\ttfamily 0123456789ABCDEFGHIJKLMNOPQRSTUVWXYZabcdefghijklmnopqrstuvwxyz}.\par
Concatenate the following 436 characters as digits $e_0,\ldots,e_{435}$;
put $\mathcal A=\sum_{j=0}^{435}e_j62^{435-j}$.\par
\begingroup\fontsize{9}{10}\selectfont\ttfamily
BxEnhvy6NipzkWmROo0GsA3uK3twWLxscLq20UTumy5Cdu461LwD6g7RDNzF8Y\par
iYDFTpB8ehAypd7d6OsOL4APamnJeNXQypUezCmW9n4oXP6XfgmiWuAWZ8ZJ7U\par
OP7gu4kVdlGZqZHULQMar1dXNeJW4i7oNVI8usp8BlovjjZbwPCovZI328fRUq\par
4q92R81s5GjWQaQfGhkq9UYMK3rRMaBmsOIBkJqc2cj2VRIZ6mjzkBIOkW9SjZ\par
FzeA6J5u4j3sBL3SdgwXzM9mHvQivu8EJ1TBTpRhHDqZIbmYgwQ8VqTAFotl6F\par
We1YC554GKrVHsA8oUGKSA6F6k5KREBSD0Tgd7bFALY84LLbpw574BQyvdoHrl\par
0Oe35uwDf15DpT6VJWZBjSbHPpPAINYTte0eUxnTFTKBlGj8Am6sb1gLldh8WZ\par
8P\par
\endgroup

\end{document}
